\section{Introduction}\label{sec:intro}
Talking head generation animates static portrait images into lifelike avatars that can speak like humans. These systems are increasingly used to create virtual presenters, hosts, and educators that can substitute for real human presence in many scenarios. They also support customized avatars that users can interact with, for instance, chatting with their favorite characters~\citep{synthesia, hedra, pika}, offering a practical tool for content creation and visual communication.

However, existing avatar generation models fail to fully replicate the feeling of real face-to-face interaction. They primarily focus on generating audio-synchronized lip or natural head motions~\citep{vasa_1, float, emo, omnihuman, magicinfinite, talkingmachines, mocha, hunyuanportrait} to deliver information accurately and naturally, rather than engaging in interactive conversations. Such one-way interaction overlooks the bidirectional nature of real-life conversations, where the continuous exchange of verbal and non-verbal signals plays a crucial role in communication. For instance, active listening behaviors such as nodding or empathic responses encourage the speaker to continue, while expressive speaking behaviors such as smiling or making eye contact contribute to a more realistic and immersive conversation.


We identify two major challenges in generating truly interactive avatars. The first challenge is the real-time processing of users’ multimodal inputs. An interactive avatar system must continuously receive and respond to the user’s multimodal cues, such as speech, head motion, and facial expression,  which requires both low inference time and minimal latency. Existing methods~\citep{infp} achieve fast inference within a motion latent space~\citep{megaportrait} but have high latency because they need the full conversation context (e.g., over 3 seconds), including future frames. The model must wait for a sufficiently long audio segment before generating motion, causing a notable delay in user interaction, as illustrated in \cref{fig:blockwise_causal_dit}. This highlights the need for a causal motion generation framework that reacts immediately to live inputs.


The second challenge is learning expressive and vibrant interactive motion. Expressiveness in human interactions is inherently difficult to define or annotate, and the lack of well-curated data makes the natural interactive behaviors hard to model. For listening behaviors in particular, we observe that most of the training data are less expressive and low-variant, often exhibiting stiff posture (See \cref{fig:vis_speaker_listener_exp}). Moreover, unlike lip synchronization, which is strongly tied with the avatar audio and therefore relatively easy to learn, reacting appropriately to user cues can correspond to a wide range of plausible motions. This ambiguity greatly increases learning difficulty, often resulting in less diverse, stiff motions, particularly in response to non-verbal cues.


To address these challenges, we present a new interactive head avatar generation framework, \textbf{Avatar Forcing}, which models the causal interaction between the user and the avatar in a learned motion latent space~\citep{lia, float}. Inspired by recent diffusion-forcing-based interactive video generative models~\citep{self_forcing, causvid, genie3}, we employ causal diffusion forcing~\citep{diffusion_forcing, causvid, streamdit} to generate interactive head motion latents while continuously processing multimodal user inputs in real-time. Unlike the previous approach~\citep{infp} that requires the future conversation context (See~\cref{fig:blockwise_causal_dit}), Avatar Forcing causally generates interactive avatar videos by efficiently reusing past information through key-value (KV) caching. This design enables real-time reactive avatar generation with minimal latency.

Moreover, following the learning-from-losing paradigm of preference optimization methods~\citep{ouyang2022training, rafailov2023dpo, wallace2024diffusiondpo}, we propose a preference optimization method to enhance the interactiveness of the avatar motion. By synthesizing under-expressive motion latents via dropping the user signals and using them as less-preferred samples,
we achieve significantly improved expressiveness and reactiveness, without the need for additional human annotations for natural interaction. As a result, Avatar Forcing can produce more natural and engaging interactive videos as illustrated in~\cref{fig:concept}.


Extensive experiments demonstrate that Avatar Forcing supports real-time interaction with a latency of roughly 500ms. Moreover, the proposed preference optimization significantly improves motion reactiveness and richness, producing more natural and engaging videos. In human evaluations, our method is preferred over 80\% against the strongest baseline, demonstrating clear advantages in naturalness, responsiveness, and overall interaction quality. 
